\subsection{A polynomial conductor for the Fermat evaluation degrees}
\label{sec:fermat-conductor}

\paragraph{P-F50 (two explicit mixed-Hessian invariants give eventual support).}
For $n\ge2$ let
\[
 \mathscr E_n=\{d\ge0:\ \exists I\in
     \mathbb C[\operatorname{Sym}^3\mathbb C^n]^{SL_n}_d,
                           \ I(F_n)\ne0\},
      \qquad F_n=\sum_i x_i^3.                         \tag{P-F50a}
\]
There are explicitly specified relatively prime positive
integers $u,v$ of polynomial size in $n$ such that
\[
          2n\langle u,v\rangle\ \subseteq\ \mathscr E_n
                         \ \subseteq\ 2n\mathbb Z_{\ge0}.
                                                        \tag{P-F50b}
\]
Here the first inclusion is realized by products of two
specified Hessian--Cayley invariants. In particular
\[
               2nr\in\mathscr E_n
             \quad\text{for all }r\ge (u-1)(v-1).
                                                        \tag{P-F50c}
\]
Unconditionally $(u-1)(v-1)=O(n^{20.72})$ with an
effective absolute implied constant. Under the supplied
Dirichlet theorem there is the explicit bound
\[
 (u-1)(v-1)<\frac14(\varphi(Q)\log Q)^4,
 \qquad Q=2n\,2^{\lceil\log_2(4n)\rceil}<16n^2.
                                                        \tag{P-F50d}
\]
Thus the bound in (P-F50d) is $O(n^8(\log n)^4)$.
In the canonical stretch variable $s$, $g_n=\gcd(n,3)$,
every $s=2g_nr$ with $r\ge(u-1)(v-1)$ has an invariant
nonzero at Fermat and hence $f_n(s)>0$.
No other stretch can be nonzero \emph{at Fermat}.
The whole invariant space can be nonzero at those other
stretches, so this last assertion is not a vanishing
claim for $f_n$.

\paragraph{A mixed-Hessian evaluation identity.}
Use the normalized Hessian $H_F$ and Cayley contraction
$\mathcal C_{D,n}$ of P-F49. Let $p$ be a prime and let
$h,j\ge0$ satisfy
\[
               nh+3j=p-1=:D,\qquad nj<p.
                                                        \tag{P-F50e}
\]
Then
\[
 J_{h,j}(F)=\mathcal C_{D,n}(H_F^h F^j)
                                                        \tag{P-F50f}
\]
is a degree-$nq$ invariant, where $q=nh+j=D-2j$, with
\[
       J_{h,j}(gF)=(\det g)^{3q}J_{h,j}(F),
                       \qquad J_{h,j}(F_n)\ne0.
                                                        \tag{P-F50g}
\]
To prove the nonzero assertion with every normalization
accounted for, first let $G=\sum_{|\alpha|=D}g_\alpha x^\alpha$
have coefficients in $\mathbb Z_{(p)}$. Grouping the
permutation terms of (P-F49f) by their exponent matrices
$L$, and then applying (P-F49h), gives
\[
 \begin{split}
 \mathcal C_{p-1,n}(G)
 &=\sum_L c_L\prod_{i=1}^n
             \left(g_{L_i}\frac{\prod_a L_{ia}!}{(p-1)!}\right)\\
 &\equiv\sum_L\prod_{i=1}^n g_{L_i}
   =[x_1^{p-1}\cdots x_n^{p-1}]G(x)^n\pmod p.
 \end{split}                                           \tag{P-F50h}
\]
Here $L_i$ is row $i$ and every row and column of $L$
sums to $p-1$. All factorial denominators are units.

At Fermat, $H_{F_n}=\prod_i x_i$. Therefore the last
coefficient in (P-F50h), for
$G=(\prod_i x_i)^h(\sum_i x_i^3)^j$, is
\[
 [x_1^D\cdots x_n^D]
       (\prod_i x_i)^{hn}(\sum_i x_i^3)^{jn}
                         =\frac{(nj)!}{(j!)^n}.        \tag{P-F50i}
\]
Indeed each exponent still required after the first
factor is $D-hn=3j$, so each of the $n$ summands must
be chosen $j$ times. The condition $nj<p$ makes this
integer nonzero modulo $p$, proving the evaluation claim.
The polynomial inserted into $\mathcal C$ has variable
degree $nh+3j=D$ and coefficient degree $nh+j=q$.
The Hessian transformation contributes determinant
exponent $2hn$, and the Cayley contraction contributes
$D$, so their sum is $3nh+3j=3q$. Its coefficient
degree is $nq$. This proves (P-F50g).

\paragraph{Two admissible mixed powers from one sufficiently large prime.}
Choose an integer $c$ with
\[
                1\le c<n,
       \qquad\gcd(c,n)=\gcd(6c+1,n)=1.                 \tag{P-F50j}
\]
Such a choice always exists. Modulo a prime divisor
$\ell=2$ or $3$ of $n$, choose any nonzero residue;
$6c+1$ is then automatically nonzero. Modulo each
$\ell\ge5$, exclude the two residues $0,-6^{-1}$;
at least three residues remain. The Chinese remainder
theorem, extended to the prime-power factors of $n$,
gives a residue modulo $n$ satisfying both coprimality
conditions. Its representative is nonzero, so lies
between one and $n-1$. Alternatively it can be found
by testing those $n-1$ integers with gcd computations.

Put
\[
 \ell=2^{\lceil\log_2(4n)\rceil},\quad Q=2n\ell,
 \quad r_0\equiv6c+1\pmod{2n},\quad 1\le r_0<2n,
 \quad a=Q/2+r_0.                                     \tag{P-F50k}
\]
The notation $\ell$ here denotes a power of two, not a
prime. It obeys $4n\le\ell<8n$, hence
$8n^2\le Q<16n^2$. The residue $r_0$ is odd and prime
to $n$. Since $\ell$ is even, $Q/2=n\ell$ is divisible
by $2n$, and consequently $a$ is odd and prime to $n$.
The only additional prime divisor of $Q$ is two.
Thus
\[
          \gcd(a,Q)=1,
                 \qquad4n^2<a<Q.                     \tag{P-F50l}
\]
Choose the least prime $p\equiv a\pmod Q$. In particular
\[
                 p>4n^2,
                 \qquad p\equiv6c+1\pmod{2n}.         \tag{P-F50m}
\]
Set
\[
 \begin{array}{lll}
 j=2c,&h=(p-1-6c)/n,&q=p-1-4c,\\
 j'=2c+n,&h'=h-3,&q'=q-2n.
 \end{array}                                           \tag{P-F50n}
\]
All these numbers are integers. To check (P-F50e) for
both pairs, note that $j'\le3n-2$ and hence
$nj'\le3n^2-2n<p$. Moreover
\[
 p-1-3j'\ge4n^2-9n+6>0\qquad(n\ge2).
\]
The last quadratic is positive because it has positive
leading coefficient and negative discriminant. Thus
$h'>0$ and $h\ge0$; both pairs meet every required
condition. The congruence modulo $2n$ makes $q,q'$ even.
Define
\[
                      u=q/2,\qquad v=q'/2=u-n.
                                                        \tag{P-F50o}
\]
They are positive since $p>4n^2$ and $c<n$.
Furthermore
\[
      u=(p-1)/2-2c\equiv c\pmod n,
          \qquad \gcd(u,v)=\gcd(u,n)=1.                \tag{P-F50p}
\]
The two instances of (P-F50f) therefore give invariants
of degrees $2nu,2nv$, both nonzero at Fermat.

\paragraph{The exact point-specific period obstruction.}
For completeness, let a homogeneous degree-$d$ $SL_n$
invariant $I$ have $I(F_n)\ne0$. The scalar center
$\mu_n\subset SL_n$ first forces $n\mid3d$.
Writing $b=3d/n$, any $g\in GL_n$ can be expressed as
a scalar times an $SL_n$ matrix, and homogeneity gives
\[
                         I(gF)=(\det g)^b I(F).
                                                        \tag{P-F50q}
\]
The diagonal matrix $\operatorname{diag}(\omega,1,\ldots,1)$,
where $\omega$ is a primitive cube root, fixes $F_n$.
Its determinant forces $3\mid b$. A transposition of
two variables also fixes $F_n$ and has determinant
$-1$, forcing $2\mid b$. Thus $6\mid b$, equivalently
$2n\mid d$. This proves the second inclusion in
(P-F50b) directly, without an assumption about the
full Fermat stabilizer. It agrees with the degree
period in B\"urgisser--Ikenmeyer, Proposition~3.24(2),
printed p.~15.

Products of the two constructed invariants stay nonzero
at $F_n$, giving the first inclusion in (P-F50b).
For relatively prime $u,v$, every integer
$t\ge(u-1)(v-1)$ belongs to $\langle u,v\rangle$.
Indeed choose $b\in\{0,\ldots,u-1\}$ with
$bv\equiv t\pmod u$. Then $t-bv$ is a multiple of $u$
and is at least $-(u-1)$, so it is nonnegative.
Setting $a=(t-bv)/u$ writes $t=au+bv$ with both
coefficients nonnegative and proves (P-F50c).

\paragraph{The unconditional and supplied-GRH bounds.}
Xylouris's Theorem~1.1, p.~66, cited in P-F49 gives
$p\le C Q^{5.18}$. Since $u,v<p/2$,
\[
       (u-1)(v-1)<p^2/4
                      \le C^2Q^{10.36}/4=O(n^{20.72}).
                                                        \tag{P-F50r}
\]
For the supplied Dirichlet theorem, apply
Lamzouri--Li--Soundararajan Corollary~1.2, p.~3,
to the reduced class $a\pmod Q$; $Q>3$ is automatic.
It gives $p\le(\varphi(Q)\log Q)^2$, and the same
inequality proves (P-F50d). The exact analytic hypotheses
and primary URLs are those checked in P-F49.

Equations (P-F50b)--(P-F50d) close the question of a
polynomial eventual-support threshold at this specific
Fermat point. They give a family of actual nonzero
cubic plethysm coefficients on all sufficiently large
stretches permitted by that point's period. They do
not claim that every invariant in those spaces is
constructed, and they give no claim of small evaluation
circuits for the Cayley contractions.

Most importantly, eventual positivity of actual values
on the subprogression $2g_n\mathbb Z$ is not a statement
about every constituent polynomial evaluated at every
positive argument after a shift. It does not settle
the strict saturation or coefficient-positivity indices
used in PH2/SH. The distinctions in P-F34 and P-F47
remain in force.

\paragraph{A full-ring support corollary on an infinite family.}
Let $r_0\ge3$ be odd, and put
\[
       k=3r_0,\qquad n=\binom{k-1}{2},\qquad
       C_n=(u-1)(v-1),
       \qquad f_n(s)=\dim\bigl(\operatorname{Sym}^{ns}
                    (\operatorname{Sym}^3\mathbb C^n)\bigr)^{SL_n}.
                                                        \tag{P-F50s}
\]
Here $u,v$ are the explicit integers constructed above for this $n$.
Since $2n=(3r_0-1)(3r_0-2)\equiv2\pmod3$, we have
$\gcd(n,3)=1$, so this is precisely the canonical stretching
grading. The top exterior-power case of (P-F32a) gives
\[
               f_n(r_0)=1,
      \qquad f_n(s)>0\quad\hbox{for every integer }
                           s\ge2C_n+r_0.              \tag{P-F50t}
\]
To verify the first assertion directly, the relevant exterior
power has degree $nr_0=\binom{k}{3}$, hence is the top exterior
power of $\bigwedge^3\mathbb C^k$. Each diagonal weight occurs
$\binom{k-1}{2}=n$ times, so that one-dimensional representation
is $(\det)^n$, with multiplicity one. Let $J$ be a nonzero
invariant in canonical degree $r_0$. For each $t\ge C_n$,
(P-F50c) supplies a nonzero invariant $I_t$ in canonical
degree $2t$. The invariant ring is a subring of the polynomial
ring in the cubic coefficients and is therefore a domain.
Thus $JI_t\ne0$, in degree $r_0+2t$.
Every even $s\ge2C_n+r_0$ is already covered by $I_{s/2}$,
and every odd such $s$ is covered by $JI_{(s-r_0)/2}$.
This proves the second assertion for all actual integer stretches,
including those outside the Fermat point's even progression.

The threshold $2C_n+r_0$ is $O(n^{20.72})$ unconditionally
and $O(n^8\log^4 n)$ under the supplied Dirichlet-GRH premise,
because $r_0=O(\sqrt n)$ and the bounds for $C_n$ were proved
above. This is a full-ring eventual-support statement on the
specified infinite subfamily of dimensions. It is still not
the strict saturation assertion that every nonzero constituent
polynomial is positive at every positive argument after a
uniform shift: actual support tests a constituent only on its
own residue progression. No PH2 or SH conclusion is added.
